\documentclass[11pt]{article}
\usepackage[margin=1in]{geometry}
\usepackage{booktabs}
\usepackage{makecell}
\usepackage{array}
\usepackage{graphicx}
\usepackage{setspace}
\usepackage{caption}
\usepackage{amsmath}
\usepackage[normalem]{ulem} % for \uline in table templates
\usepackage{dsfont} % for \mathds in parameter labels
\captionsetup{font=small,labelfont=bf}

\title{Startup Definition Robustness: Size-Based Cutoffs}
\date{\today}

\begin{document}
\maketitle
\onehalfspacing

\section*{Overview}
This note documents the size-based startup definition (firm headcount at 2019H2) and reports the corresponding regression tables and cutoff-sweep figures for user productivity and firm scaling outcomes. Startup is defined as $\text{size}_{2019H2} \leq c$, with the main cutoff at 150 employees and robustness across \{50, 100, 150, 200, 500\}.

\section*{Main Tables (size cutoff = 150 employees)}

\subsection*{User Productivity}
\input{../../results/cleaned/tex/startup_size_cutoff_total_contributions_q100.tex}
\input{../../results/cleaned/tex/startup_size_cutoff_restricted_contributions_q100.tex}

\subsection*{Firm Scaling}
\input{../../results/cleaned/tex/startup_size_cutoff_growth_rate_we.tex}
\input{../../results/cleaned/tex/startup_size_cutoff_join_rate_we.tex}
\input{../../results/cleaned/tex/startup_size_cutoff_leave_rate_we.tex}

\end{document}
